\documentclass[11pt,a4paper]{report}
\usepackage{marvosym}

\assignment{1}
\group{G095}
\students{Couvreur Matthieu}{Benbada Abdellatif}

\begin{document}

\maketitle

Answer to the questions by adding text in the boxes. 
You may answer in either \textbf{French or English}. 
Do not modify anything else in the template.  
The size of the boxes indicate the place you \textbf{can} use, but you \textbf{need not} use all of it (it is not an indication of any expected length of answer). 
\textbf{Be as concise as possible! A good short answer is better than a lot of nonsense!}
%\bigskip

\section{Understanding a CP model (2 points)}

\begin{enumerate}
    \item Answer correctly to the questions 1-4. \textbf{(0.5pt)}
    You can help yourself by looking
    at \url{https://www.pycsp.org/documentation/constraints/} to find the definition of the different constraints.
    \item Answer correctly to the question 5. \textbf{(0.5pt)}
    \item Add the needed constraint to the model and submit it on INGInious. \textbf{(1pt)}
\end{enumerate}

\newpage
\section{Minesweeper Model (8 points)}

\begin{enumerate}
    \item Implement your CP model in the file \texttt{minesweeper.py} and submit it on INGInious. (\textbf{4 pts})
    \item Describe your CP model: First describe the variables of your CP model. 
    For each set of variables, explain what they represent, give their domains and how many variables are in this set. 
    Then, list the constraints of your CP model. 
    For each constraint explain what it enforces, and on which variables it is applied. (\textbf{4 pts})
\end{enumerate}

\begin{answers}[10cm]
    \item Variables of the model : \\
    In the backtracking algorithm, we need an assignment variable, the name is the same in the program. Here we create an Array $assignements$ of the size of the map with 2 possible domains : {0 : not a mine, 1 : a mine}. $neighbours$ represent the neighbours of the studied cell around. In the backtracking method we use SELECT-UNASSIGNED-VARIABLE, meaning which next value are we studying. One solution is to find the value with the most constraints, here its the one with the biggest number. Thats what $priority\_queue$ is used for, we save the value of the cell and its coordinates then sort by the value. 
    
    \item Constraints of the CP model : The constraints are applied to the grid configuration, we will say if the cell should be a mine (value = 1) or not (value = 0). In order to put constraints by using the clues, we first need to check if the cell has a clue to give ($local\_clue!=-1$ means that its not an empty cell). As stated in the rules of minesweeper, the number written in the studied cell must be equal to the number of mines around it. This is why we are making a loop to get every neighbour of the cell so we can make the 1st constraint. The constraint is written here with $Sum(neighbours) == local\_clue$ (with $local\_clue$ the clue of the studied cell) that we implement with the $satisfy$ function given by pycsp3 thanks to the bibliography. Then if we have a clue to study (the value is not -1), we already know that the cell is not a mine, we can already put it as a constraint $assignements[i][j] == 0$. 
    
    \item Solving : 
    After filling $assignements$, we check if its solved or not thanks to pycps3 functions, if not it will return None. If the problem is solved, we go through the $assignements$ and take each cell who is a mine (value = 1) in a form of a tuple just as asked in the subject.
    % Your answer here
\end{answers}




\newpage
\section{Atom Placement Problem()}
\begin{enumerate}
\item Formulate the atom placement problem as a Local Search problem (problem, cost function, feasible solutions, optimal solutions). \textbf{(1 pt)}
\end{enumerate}

\begin{answers}[6cm]
% Your answer here
\end{answers}


\begin{enumerate}
\setcounter{enumi}{1}
    \item You are given a template on Moodle: \textit{atomplacement.py}. Implement your own extension of the \textit{Problem} class from \texttt{aima-python3}. Implement the \texttt{maxvalue} and \texttt{randomized maxvalue} strategies. To do so,
	you can get inspiration from the \textit{randomwalk} function in \textit{search.py}. Your program will be evaluated on 15 instances (during 1 minute) of which 5 are hidden. We expect you to solve at least 12 out the 15.
    \begin{enumerate}
        \item \texttt{maxvalue} chooses the best node (i.e., the node with maximal
        value) in the neighborhood, even if it degrades the quality of the current solution. 
        The \texttt{maxvalue} strategy should be defined in a function called
        \textit{maxvalue} with the following signature: \\
        \texttt{maxvalue(problem,limit=100)}. \textbf{(2.5 pts)}
    \end{enumerate}
\end{enumerate}

\begin{answers}[2.5 cm]
% IN CASE YOU WANT TO SPECIFY ANY COMMENTS ABOUT YOUR IMPLEMENTATION
\end{answers}



\begin{enumerate}
\setcounter{enumi}{1}
\begin{enumerate}
\setcounter{enumii}{1}
    \item \texttt{randomized maxvalue} chooses the next node randomly among
        the 5~best neighbors (again, even if it degrades the quality of the current solution).
        The \texttt{randomized maxvalue} strategy should be defined in a function called
        \textit{randomized\_maxvalue} with the following signature: \\
        \texttt{randomized\_maxvalue(problem,limit=100)}. \textbf{(2.5 pts)}
\end{enumerate}
\end{enumerate}

\begin{answers}[2.5 cm]
% IN CASE YOU WANT TO SPECIFY ANY COMMENTS ABOUT YOUR IMPLEMENTATION
\end{answers}



\begin{enumerate}
\setcounter{enumi}{2}
\item Compare the 2 strategies implemented in the previous question and \texttt{randomwalk} defined in \textit{search.py} on the given vertex cover instances. 
    Run the strategies during 100 steps, and report, in a table, the computation
    time, the value of the best solution and the number of steps needed to reach the best result. For the
    randomized max value and the random walk, each instance should be tested 10~times to reduce
    the effects of the randomness on the result. When multiple runs of the same instance
    are executed, report the mean of the quantities. \textbf{(3 pts)}
\end{enumerate}

\begin{answers}[7cm]
% Your answer here
\begin{center}
\begin{tabular}{||l||l|l|l||l|l|l||l|l|l||}
\hline
\multirow{3}{*}{Inst.} & \multicolumn{3}{c||}{Max value} & \multicolumn{3}{c||}{Random max value} & \multicolumn{3}{c||}{Random walk} \\
\cline{2-10}
\cline{2-10}
 & T(s) & Val & NS & T(s) & Val & NS & T(s) & Val & NS\\
\hline
i01 & & & & & & & & &\\
\hline
i02 & & & & & & & & &\\
\hline
i03 & & & & & & & & &\\
\hline
i04 & & & & & & & & &\\
\hline
i05 & & & & & & & & &\\
\hline
i06 & & & & & & & & &\\
\hline
i07 & & & & & & & & &\\
\hline
i08 & & & & & & & & &\\
\hline
i09 & & & & & & & & &\\
\hline
i10 & & & & & & & & &\\
\hline
\end{tabular}
\end{center}
\textbf{NS}: Number of steps — \textbf{T}: Time — \textbf{Val}: Value of the best solution
\end{answers}



\begin{enumerate}
\setcounter{enumi}{3}
    \item \textbf{(4 pts)} Answer the following questions:
    \begin{enumerate}
        \item In one sentence, what is the best strategy? \textbf{(0.5 pt)}
    \end{enumerate}
\end{enumerate}

\begin{answers}[3.5cm]
% Your answer here
\end{answers}


\newpage
\begin{enumerate}
\setcounter{enumi}{3}
\begin{enumerate}
\setcounter{enumii}{1}
    \item Why do you think the best strategy beats the other ones? What conclusions can be drawn on the atom placement problem ? \textbf{(1.5 pt)}
\end{enumerate}
\end{enumerate}

\begin{answers}[7.5cm]
% Your answer here
\end{answers}



\begin{enumerate}
\setcounter{enumi}{3}
\begin{enumerate}
\setcounter{enumii}{2}
    \item What are the limitations of each strategy in terms of diversification 
    and intensification? \textbf{(1 pt)}
\end{enumerate}
\end{enumerate}

\begin{answers}[5cm]
% Your answer here
\end{answers}



\begin{enumerate}
\setcounter{enumi}{3}
\begin{enumerate}
\setcounter{enumii}{2}
    \item What is the behavior of the different techniques when they fall 
    in a local optimum? \textbf{(1 pt)}
\end{enumerate}
\end{enumerate}

\begin{answers}[5cm]
% Your answer here
\end{answers}


\end{document}
